\subsection{Experimental Setup and Assessment Methodology}

In this section, the experimental setup for evaluating the proposed autonomous antenna
alignment mechanism, metrics to measure performance, and comparison protocol with other
algorithms will be discussed. The aim is to make sure that all algorithms are
evaluated under identical initial state, actuator constraints, and RSSI measurement
procedure.

\subsubsection{Experimental Environment}

The experimental setting consists of two main environments that simulate different
propagation environments:

\begin{itemize}
    \item \textbf{Indoor setting:} Lab or corridor-like settings where reflective
    surfaces like walls and metallic items exist. Multipath propagation and RSSI
    variation become more pronounced in such settings.
    \item \textbf{Outdoor setting:} Open or semi-open environments where the
    distance between the transmitter and receiver varies from 10 to 30~m in practice.
    This setting exhibits line-of-sight behavior with some obstacles present.
\end{itemize}

The experiments are performed in the presence of the following assumptions:
\begin{itemize}
    \item The orientation and the transmission power of the transmitter are held
    constant during an experiment.
    \item The orientation of the receiver is started with a deliberate initial misalignment
    chosen from the same predetermined set for all compared algorithms.
    \item In order to compare the performance of the algorithms practically, static tests
    are made under a disturbance-free environment.
    \item Repeated experiments are carried out for different practical situations.
\end{itemize}

\subsubsection{Evaluation Metrics}

In order to allow for a quantitative comparison of the proposed RL-based alignment scheme with
baselines approaches, the following performance measures are collected:

\textbf{Received Signal Strength Indicator (RSSI):} 
This measure is considered to be a main metric of evaluating the quality of the communication
link. The RSSI value used for analysis corresponds to the averaged one measured within
each decision step as an average over 20 packets with 10~ms interval between them. 
This parameter is considered to be meaningful only if it provides additional gain on
the received signal strength level with respect to the initial one observed at the
wireless receiver side. In addition, the RSSI is directly used by the RL agent as an
observation and is thus a natural choice for a fundamental comparison basis.

\begin{equation}
RSSI_{\text{avg}}(k) = \frac{1}{N}\sum_{i=1}^{N} RSSI_i(k), \qquad N = 20
\end{equation}

\textbf{Convergence Time ($T_{conv}$):} 
This parameter indicates the amount of time required to achieve a steady state within the practical
loop. Within the deployment firmware implementation, the convergence criterion is met if the absolute
difference of the averaged RSSI does not exceed 0.8~dB for four subsequent control cycles, or the maximum
number of decision steps is reached. This criterion is used consistently during the comparison of all methods
in practical conditions to ensure comparability of the reported values.
\begin{equation}
\left| RSSI_{\text{avg}}(k) - RSSI_{\text{avg}}(k-1) \right| < 0.8~\text{dB},
\quad \forall k \in [k_c, k_c+3]
\end{equation}
\begin{equation}
T_{conv} = t(k_c)
\end{equation}

0.8~dB was chosen as a practical dead band for further hardware implementation. The value is high enough to
consider the achieved RSSI steady enough to be considered as a high gain value, yet low enough not to be
affected by fluctuations in packet-level measurements and actuator dynamics.

\textbf{Movement Count:} 
Total number of executed movements of the commanded pan and tilt operations before convergence.
Number of movements is chosen to estimate the actuation effort due to the inverse relation between motions
and wear of the mechanical elements, energy consumption, and time spent on alignment.
\begin{equation}
M = N_{\text{pan}} + N_{\text{tilt}}
\end{equation}

\textbf{Success Rate:} 
The success rate is defined as the percentage of trials in which the method succeeds in converging
to a high RSSI solution. This metric serves to differentiate fast algorithms that often yield
erroneous solutions from those consistently achieving correct solutions.
\begin{equation}
SR(\%) = \frac{N_{\text{success}}}{N_{\text{trials}}} \times 100
\end{equation}

\textbf{RSSI Variability:} 
Spread of final RSSI values after each repeated trial reported in the results section by using
the standard deviation of RSSI. The smaller the variance, the more reliable the method,
i.e., it is less prone to variations due to changes in environment and other factors.
\begin{equation}
\sigma_{\text{RSSI}} = \sqrt{\frac{1}{N_{\text{trials}}}\sum_{j=1}^{N_{\text{trials}}}
\left(RSSI_{\text{final},j} - \overline{RSSI}_{\text{final}}\right)^2}
\end{equation}

\subsubsection{Experimental Protocol}

Every test follows a protocol designed for reproducibility:

\begin{itemize}
    \item The algorithm starts from one of the defined pairs of an azimuth and tilt 
    values corresponding to initial misalignment states.
    \item The algorithm is run starting from the selected state in a consistent 
    simulation environment.
    \item Runtime parameters are consistent between the algorithms including RSSI 
    sampling rate, actuator limitations, and time delay before settling down.
    \item Every control loop logs a timestamp, averaged RSSI readings, 
    orientation indices, performed actions, and whether alignment occurred.
    \item Each practical situation is simulated 10 times using the selected 
    method for every compared approach.
    \item For a comparison, the presented reinforcement learning method is tested against
    exhaustive search and hill-climbing based on RSSI values.
\end{itemize}

\subsubsection{Statistical Analysis}

For statistical indicators, the average value over several repetitions serves as a 
benchmark for the performance comparison, whereas standard deviations can provide 
additional information on consistency of results. As all tests are done on matching 
starting states with equal runtime configurations, the performed analysis allows to compare 
algorithms in practice.